\begin{definition}
Let $\Omega\subset\mathbb{C}$ be open and let $1\le K<\infty$.  A mapping
$\chi\colon\mathbb{C}\to\mathbb{C}$ is a \emph{$K$-quasiconformal
homeomorphism of $\Omega$} if
\begin{enumerate}
\item[(a)] $\chi$ is $K$-quasiregular on $\Omega$ in the sense of
\noderef{QuasiregularOn};
\item[(b)] $\chi$ is injective on $\Omega$ and continuous on $\Omega$;
\item[(c)] the image $\chi(\Omega)$ is open; and
\item[(d)] there is a map $\psi$ with $\psi(\chi(z))=z$ for every $z\in\Omega$
which is continuous on $\chi(\Omega)$.
\end{enumerate}

Conditions (b)--(d) say that $\chi$ restricted to $\Omega$ is a homeomorphism
of $\Omega$ onto the open set $\chi(\Omega)$, with inverse the restriction of
$\psi$: injectivity and (d) give $\psi\circ\chi=\mathrm{id}$ on $\Omega$, and
for $w=\chi(z)\in\chi(\Omega)$ one gets $\chi(\psi(w))=\chi(z)=w$, so $\psi$
restricted to $\chi(\Omega)$ is a two-sided inverse; both maps are continuous
on the relevant sets.

This is the notion of quasiconformal homeomorphism used throughout the paper,
in particular for the factor $\chi$ of the factorization (15) and in the
standing assumptions of Sections~3 and~5.  A quasiconformal homeomorphism is
by definition a quasiregular mapping which is in addition a homeomorphism
onto its image.
\end{definition}
